\documentclass[10pt,a4paper]{amsart}
\usepackage[margin=19mm]{geometry}
\usepackage{amsmath,amssymb,mathtools}
\usepackage[scr=rsfs,cal=euler]{mathalfa}
\usepackage{xfrac}
\usepackage[unicode,hidelinks,bookmarks=false]{hyperref}
\newcommand{\ie}{i.e.\ }
\newcommand{\cf}{cf.\ }
\newcommand{\resp}{respectively\ }
\newcommand{\Subsection}{\subsection}
\providecommand{\nobreakdash}{\nobreak\hbox{-}}
\setlength{\emergencystretch}{2em}
\multlinegap0pt
\usepackage{amssymb}
\usepackage[shortlabels]{enumitem}
\usepackage{xcolor}
\newtheorem{lemma}{Lemma}
\newcommand{\R}{\mathbb R}
\title{Fokker-Planck Analysis and Invariant Laws for a Continuous-Time Stochastic Model of Adam-Type Dynamics}
\author{Kaj Nystr{\"o}m}
\date{Source: arXiv:2604.00840v1 (CC BY 4.0)}
\begin{document}
\maketitle

\noindent\emph{Source and licence.} Fokker-Planck Analysis and Invariant Laws for a Continuous-Time Stochastic Model of Adam-Type Dynamics. Version arXiv:2604.00840v1 (CC BY 4.0), distributed by arXiv under the Creative Commons Attribution 4.0 International licence, http://creativecommons.org/licenses/by/4.0/, as recorded in the arXiv licence metadata for this exact version and shown on the abstract page https://arxiv.org/abs/2604.00840v1. Full licence notice: \texttt{licenses/arxiv-2604.00840-CC-BY-4.0.txt}.\par\smallskip

\noindent\textit{Editorial scope.} Counted unit: the article's lemma lem:power-Lyap (source0.tex lines 4561--4579). The article's complete original proof of it is delivered with it (source0.tex lines 4580--4683, one \texttt{proof} environment of 104 lines). No mathematics is written, repaired or re-derived: the statement and the proof are the article's own lines, in the article's own notation, and no retained line is substituted or re-flowed.  Retained uncounted as the source-local context the proof needs: lines 1798--1802.  Nothing was omitted from any retained span, so the package carries no omission record.  No retained line contains a citation, so the wrapper carries no bibliography and no bibliography entry.  No retained line contains a figure environment, an image-inclusion command or drawing code, so no image is delivered and the package declares no asset dependency.  Editorial formatting only, in addition: an AMS \texttt{amsart} preamble that restates the article's own declarations and macros used by the retained lines, namely the environments \texttt{lemma} on separate counters, together with the standard packages those lines need.  The retained environments print in this document as Lemma 1 in order of appearance, whereas the article prints the same items with the numbering of its own full document.  The complete native source of the article as distributed in the arXiv e-print for this exact version is bundled unchanged as \texttt{sources/197614/source0.tex}.
\begin{equation}\label{V}
V(x,z,y)=V_\upsilon(x,z,y)
= \theta\bigl(f(x)-f_\ast\bigr) + \frac\alpha2\|z\|^2 - \beta\,x\cdot z + \delta \|y\|_{1,\upsilon}
\qquad \theta,\alpha,\beta,\delta, \upsilon>0,
\end{equation}

\begin{lemma}\label{lem:power-Lyap} Assume that $f$ satisfies \textnormal{(A1)} and \textnormal{(A2)}. Let $V$ be given by \eqref{V} with $\alpha=1$. Assume that there exist
$\lambda',K'>0$, a compact set $C\subset\R^{3m}$, and constants $c_1,c_2>0$ such that, for all $(x,z,y)\in \mathbb{R}^m \times \mathbb{R}^m \times \mathbb{R}^m$, and for all $\epsilon\in [0,\upsilon]$, $\upsilon\in (0,1)$,
\[
{\mathcal L_\epsilon}V_\upsilon \;\le\; -\lambda' V_\upsilon + K'\,\mathbf 1_{C},
\qquad
V_\upsilon(x,z,y)\;\ge\; \hat c_1\bigl(\|x\|^2+\|z\|^2+\|y\|_{1,\upsilon}\bigr)-\hat c_2.
\]
Then there exists $\eta>0$ small and
constants $\lambda^{''}, K^{''}>0$, all independent of $\epsilon$ and $\upsilon$, such that
\begin{equation}\label{eq:power-drift}
{\mathcal L_\epsilon}\,\Phi(V_\upsilon) \;\le\; -\,\lambda^{''}\,\Phi(V_\upsilon)\;+\; K^{''}\,\mathbf 1_{C},
\qquad
\Phi(s):=(1+s)^{1+\eta}.
\end{equation}
Consequently,
\[
\sup_{t\ge0}\,\mathbb E_{\mu_0}\bigl[\Phi\!\bigl(V_\upsilon(X_t)\bigr)\bigr]\;<\;\infty.
\]
\end{lemma}

\begin{proof}
Let $\Phi(s)=(1+s)^{1+\eta}$ with $\eta\in(0,1)$ to be chosen small.
We compute ${\mathcal L}_\epsilon\Phi(V_\upsilon)$ using the chain rule for generators.
Since the diffusion acts only in the $z$–variables, we obtain
\[
{\mathcal L}_\epsilon\,\Phi(V_\upsilon)
=\Phi'(V_\upsilon)\,{\mathcal L}_\epsilon V_\upsilon
+\frac{a^2\sigma^2}{2}\Phi''(V_\upsilon)
\sum_{i=1}^m(\partial_{z_i}V_\upsilon)^2 .
\]
Here
\[
\Phi'(s)=(1+\eta)(1+s)^{\eta},
\qquad
\Phi''(s)=\eta(1+\eta)(1+s)^{\eta-1}.
\]

\medskip

\noindent
Since $\alpha=1$, we have
\[
V_\upsilon(x,z,y)
= \theta(f(x)-f_\ast)
+ \frac12\|z\|^2
- \beta\, x\cdot z
+ \delta\|y\|_{1,\upsilon},
\]
and we estimate
\[
\sum_{i=1}^m(\partial_{z_i}V_\upsilon)^2
= \|z-\beta x\|^2
\le 2\|z\|^2 + 2\beta^2\|x\|^2.
\]
By the coercive lower bound on $V_\upsilon$ we see that there exist $C_0>0$, $C_1>0$, independent of $\epsilon$ and $\upsilon$,
such that
\[
\|z\|^2+\|x\|^2
\le C_0(1+V_\upsilon)\implies
\sum_{i=1}^m(\partial_{z_i}V_\upsilon)^2
\le C_1(1+V_\upsilon).
\]
Therefore,
\[
\frac{a^2\sigma^2}{2}\Phi''(V_\upsilon)
\sum_{i=1}^m(\partial_{z_i}V_\upsilon)^2
\le C_2(1+V_\upsilon)^{\eta},
\]
for some constant $C_2>0$ independent of $\epsilon,\upsilon$. By assumption,
\[
{\mathcal L}_\epsilon V_\upsilon
\le -\lambda' V_\upsilon + K'\mathbf 1_{C},
\]
uniformly in $\epsilon\in[0,\upsilon]$, $\upsilon\in(0,1)$.
Hence
\[
{\mathcal L}_\epsilon\Phi(V_\upsilon)
\le (1+\eta)(1+V_\upsilon)^{\eta}
\bigl(-\lambda' V_\upsilon + K'\mathbf 1_{C}\bigr)
+ C_2(1+V_\upsilon)^{\eta}.
\]
Enlarge $C$ if necessary so that $V_\upsilon\ge v_0>0$ on $C^c$.
Then on $C^c$,
\[
(1+V_\upsilon)^{\eta}V_\upsilon
\ge \frac{v_0}{1+v_0}(1+V_\upsilon)^{1+\eta}
=: c_\star\,\Phi(V_\upsilon)\implies
(1+V_\upsilon)^{\eta}
\le \frac{1}{1+v_0}\Phi(V_\upsilon).
\]
Therefore, on $C^c$,
\[
{\mathcal L}_\epsilon\Phi(V_\upsilon)
\le -\lambda'(1+\eta)c_\star\,\Phi(V_\upsilon)
+ \frac{C_2}{1+v_0}\Phi(V_\upsilon).
\]
Choosing $\eta>0$ sufficiently small and enlarging $C$ if necessary,
we obtain
\[
{\mathcal L}_\epsilon\Phi(V_\upsilon)
\le -\lambda^{''}\Phi(V_\upsilon)
\qquad \text{on } {C}^c,
\]
for some $\lambda^{''}>0$ independent of $\epsilon,\upsilon$. Since $C$ is compact and ${\mathcal L}_\epsilon\Phi(V_\upsilon)$
is continuous, there exists $K^{''}>0$,
independent of $\epsilon,\upsilon$, such that
\[
{\mathcal L}_\epsilon\Phi(V_\upsilon)
\le K^{''}\mathbf 1_{C}.
\]
Combining the estimates on $C$ and ${C}^c$ yields
\[
{\mathcal L}_\epsilon\Phi(V_\upsilon)
\le -\lambda^{''}\Phi(V_\upsilon)
+ K^{''}\mathbf 1_{C},
\]
which proves \eqref{eq:power-drift}. Applying Dynkin's formula and Grönwall's inequality then gives
\[
\sup_{t\ge0}
\mathbb E_{\mu_0}\bigl[\Phi(V_\upsilon(X_t))\bigr]
<\infty,
\]
uniformly in $\epsilon$ and $\upsilon$. This completes the proof.
\end{proof}

\end{document}
